\begin{answer}

\textbf{(b) How the evaluation was conducted}

For each of the 30 instructions in \texttt{alpaca\_eval\_first\_30.jsonl} I
sampled one response from model E and one from model F, passing the instruction
as a single user turn with no additional prompting. I then built a judge prompt
containing the instruction and both responses, queried \textbf{model Z} with it,
and parsed a preference out of Z's reply. Win rates are the fraction of the 30
problems on which each model was preferred.

I ran this whole procedure \textbf{twice}: once with model E's response placed
first in the judge prompt, and once with model F's placed first, keeping each
label attached to its true model. The reason is given in part (c) --- the judge
turns out to be strongly sensitive to presentation order, so a single order
does not yield an interpretable win rate. Every model response and judge reply
was written to a \texttt{.jsonl} file as it was produced, which is what parts (c)
and (d) analyse.

\textbf{The judge prompt.} Model E first (the position-swapped variant is
identical with the two response blocks, the two listed tags, and the worked
example exchanged):

\begin{lstlisting}[language={},numbers=none,basicstyle=\footnotesize\ttfamily,frame=single]
You will be given two responses from different models to the same user query.
Compare them objectively and comprehensively, then select the better one.

Query: {query}

response_E: {response_E}

response_F: {response_F}

Now report your result. The LAST line of your response must be your final
choice, formatted exactly like this:

#### <MODEL_E_BETTER>

or

#### <MODEL_F_BETTER>

For example, if you think responce_E is better, the last line must read
exactly:

#### <MODEL_E_BETTER>
\end{lstlisting}

\textbf{How scores were extracted.} The judge writes free-form prose, so the
prompt above requires it to end with one of two fixed sentinel tags on its final
line. Extraction applies the regular expression
\texttt{\#\#\#\#\textvisiblespace(<MODEL\_[EF]\_BETTER>)} and takes the
\textbf{last} match in the reply. Taking the last rather than the first match
matters: the judge routinely restates both options while reasoning --- writing
things like ``\texttt{\#\#\#\# <MODEL\_E\_BETTER> would be wrong here}''
mid-paragraph --- so the first match often captures a mention rather than the
verdict. If no tag matches at all the function returns
\texttt{<NO\_PREFERENCE\_FOUND>} and the comparison is excluded from the win
rate. In practice this never fired: all 60 judge replies across both runs
contained an extractable tag.

\textbf{Results, and which model was better.}

\begin{figure}[h]
\centering
\includegraphics[width=0.85\textwidth]{judge_win_rates.png}
\caption{Judge win rates for models E and F, shown separately for each
presentation order and pooled over both. The dashed line marks the 50\% a
position-indifferent judge would produce.}
\end{figure}

\begin{center}
\begin{tabular}{lcc}
\toprule
Presentation order & Model E win rate & Model F win rate \\
\midrule
Model E shown first & 30/30 \;(100\%) & 0/30 \;(0\%) \\
Model F shown first & 11/30 \;(37\%)  & 19/30 \;(63\%) \\
\midrule
Pooled (counterbalanced) & \textbf{41/60 \;(68\%)} & 19/60 \;(32\%) \\
\bottomrule
\end{tabular}
\end{center}

\textbf{Model E is judged the better model}, winning 68\% of the 60
counterbalanced comparisons. That headline number should be read with the two
caveats established in parts (c) and (d): it depends on averaging over
presentation orders, since either order on its own gives a wildly different
answer (100\% versus 37\%), and model E is the longer response in every single
comparison, so ``E is better'' cannot be separated from ``the judge prefers
length'' using this data.

\vspace{1em}
\textbf{(c) Comparing my own preferences against the judge's}

I inspected the first 5 AlpacaEval problems together with model E's response,
model F's response, and the judge's full response for each.

\textbf{Agreement with the judge.} Of the 5 problems I inspected, I
\textbf{agreed} with the judge on \textbf{2} and \textbf{disagreed} on
\textbf{2}. On the remaining problem I had no clear preference between the two
responses --- both were accurate and well organised, and I could not honestly
say either was more useful --- so I have not scored it either way. Notably, the
judge expressed no such indecision anywhere: it returned a definite winner on
all 30 problems and never once declined to choose.

\textbf{An example where I disagree.} This is problem index 0 in
\texttt{llm\_judge\_results.jsonl}; the excerpts below are shortened, and the
full texts are in that file.

\textbf{AlpacaEval query:}

\begin{lstlisting}[language={},numbers=none,basicstyle=\footnotesize\ttfamily,frame=single]
What are the names of some famous actors that started their careers on Broadway?
\end{lstlisting}

\textbf{Model E's response} (1120 output tokens, 15 actors named):

\begin{lstlisting}[language={},numbers=none,basicstyle=\footnotesize\ttfamily,frame=single]
Broadway has long been a hallowed ground for aspiring actors, a proving ground
where talent is honed and stars are often born. Many of the most recognizable
and beloved actors in film and television today owe their early careers to the
bright lights of the Great White Way. Here are some of the most famous actors
whose journeys began on the Broadway stage:

**Musical Theatre Maestros and Dramatic Powerhouses:**

*   **Meryl Streep:** While she's now an icon of cinema, Meryl Streep's
    theatrical roots are deep. She made her Broadway debut in the 1975
    production of *Trelawney of the 'Wells'* and quickly established herself as
    a formidable stage actress with roles in *27 Wagons Full of Cotton* and
    *Secret Service*. Her early theatrical work showcased the incredible range
    and emotional depth that would later define her legendary film career.

    [... 14 further actors, each with a similar paragraph ...]
\end{lstlisting}

\textbf{Model F's response} (879 output tokens, 20 actors named):

\begin{lstlisting}[language={},numbers=none,basicstyle=\footnotesize\ttfamily,frame=single]
Broadway has been a launchpad for countless incredibly talented actors who went
on to achieve massive fame in film and television. Here are some of the most
famous ones, categorized for clarity:

**From Musicals to Stardom:**

*   **Audra McDonald:** A true Broadway legend with six Tony Awards, she's also
    a celebrated television and film actress.

    [... 19 further actors, each in one or two lines ...]
\end{lstlisting}

\textbf{Judge's response} (conclusion; it chose model E):

\begin{lstlisting}[language={},numbers=none,basicstyle=\footnotesize\ttfamily,frame=single]
Model E's strength lies in its detailed descriptions for each actor. It goes
beyond just naming them and provides a brief narrative about their Broadway
beginnings [...]

Model F includes a slightly wider range of actors, bringing in names like Lea
Michele, Matthew Morrison, Sara Bareilles, Anna Kendrick, Al Pacino [...]

While Model F offers a broader selection of names, Model E provides a richer,
more informative experience for the names it does include. The depth of detail
in Model E makes it slightly more valuable for someone looking to understand
*how* these actors started on Broadway.

#### <MODEL_E_BETTER>
\end{lstlisting}

\textbf{Why I disagree.} The query asks for \emph{names}. Model F returns 20 of
them in 3{,}697 characters; model E returns 15 in 5{,}429. F therefore delivers
33\% more of the requested content in 32\% less text, and reaches its first name
after 231 characters of preamble against E's 419. On the task as actually posed,
F dominates E on both axes that matter.

What makes this case diagnostic is that the judge \emph{identifies} the tradeoff
and then resolves it the wrong way: it writes that ``Model F offers a broader
selection of names'' before preferring E for being ``richer, more informative.''
It is trading away five of the requested items in exchange for prose depth the
user never asked for, and justifying this by reference to a goal it invented
itself --- ``someone looking to understand \emph{how} these actors started on
Broadway'' --- which is a rewriting of a request for a list into a request for an
essay. This is the mechanism by which a length-biased judge inflates a verbose
model's win rate: not by miscounting content, but by silently substituting
``more elaborate'' for ``more useful.''

\textbf{Is anything consistently different between model E and model F?}

Yes, and it is completely consistent: \textbf{model E's response is longer than
model F's on every single one of the 60 comparisons} I collected (30 problems
$\times$ 2 presentation orders). Measured in output tokens, model E averages
1547 tokens against model F's 1001 in the first run and 1533 against 1025 in the
second --- a ratio of roughly $1.5\times$ in both. Qualitatively, model E's
answers are expansive: they open with a framing paragraph, use section headings
and bulleted lists, and elaborate each point, whereas model F's answers tend to
state the requested content and stop.

This has an important consequence for the evaluation: because model E is
\emph{always} the longer response, ``the judge chose model E'' and ``the judge
chose the longer response'' are the same event in this dataset. The two
explanations are perfectly confounded and no amount of analysis of these 60
comparisons alone can separate them.

\textbf{Is the LLM judge doing a good job?}

No. The judge has a large, measurable \textbf{position bias}, visible in the
table in part (b): with model E shown first it picks E every single time, and
with model F shown first it picks E only 37\% of the time. Aggregating over both
orders, the judge selected \textbf{whichever response was shown first in 49 of 60
comparisons (82\%)}, against the 50\% a position-indifferent judge would
produce. A judge whose verdict depends this strongly on an arbitrary
presentation choice is not measuring response quality reliably, and the 100\%
win rate obtained from a single fixed order is substantially an artifact of that
ordering rather than a finding about the models.

\vspace{1em}
\textbf{(d) Response length versus judge preference}

\begin{figure}[h]
\centering
\includegraphics[width=0.85\textwidth]{judge_length_bias.png}
\caption{Distribution of response lengths, in output tokens, for the responses
the judge preferred and the responses it did not, pooled over both presentation
orders (60 comparisons). Bars are drawn side by side within each bin rather than
overlaid, so neither group is obscured.}
\end{figure}

\textbf{Note on the data used.} The problem statement anticipates 30 datapoints
per histogram, from a single pass over the 30 problems. I have pooled both
presentation orders instead, giving 60 per histogram, because a single fixed
order confounds length with position (see part (c)); the pooled version is the
one that isolates length. For reference, the single-order version is far more
dramatic --- with model E always shown first the judge picked model E every
time, so the ``preferred'' histogram would simply be model E's length
distribution (mean 1547 tokens) and the ``not preferred'' histogram model F's
(mean 1001), separating almost perfectly with no overlap. No comparison had to
be discarded: all 60 judge responses yielded an extractable preference.

\textbf{Trends.} In the pooled data the preferred responses are longer on
average --- mean 1403 tokens versus 1150, a median of 1389 versus 1065 --- and
the judge picked the longer of the two responses in \textbf{41 of 60
comparisons (68\%)}. So a length preference is visible, but the two
distributions overlap heavily; length clearly does not determine the verdict on
its own. The contrast between this and the single-order result is itself the
main finding: the apparently overwhelming length effect under a fixed
presentation order shrinks to a moderate one once position is counterbalanced,
which means most of the original effect was position, not length.

\textbf{Do I trust this evaluation as a proxy for user preference?}

No, for three reasons that compound.

First, the judge is not order-invariant. It favours the first-presented response
82\% of the time, so a quantity as arbitrary as which response I happened to
paste first moves the headline win rate from 100\% to 37\%. Any single-order win
rate --- including the 100\% I originally measured --- is therefore not a
property of the models being compared.

Second, even after counterbalancing removes the position effect, the residual
signal cannot be interpreted as quality. Model E is the longer response in all
60 comparisons, so its remaining 68\% win rate is equally consistent with ``E is
genuinely better'' and with ``the judge rewards length.'' Distinguishing these
would require comparisons in which the terser response is sometimes the longer
one --- for instance by varying verbosity independently of model identity, or by
truncating and padding responses to equalise length.

Third, none of this is evidence about human preference at all. Longer, heavily
formatted answers are exactly the kind of output an LLM judge rewards, but real
users frequently prefer a direct answer to a well-organised essay, particularly
for the short factual requests that make up much of AlpacaEval. Validating this
pipeline would require collecting actual human preferences on the same 30
problems and measuring how often the judge agrees --- which is what part (c)
begins to do on a very small scale, and my own disagreements there are a first
indication that judge preference and human preference come apart.

\end{answer}
